\chapter{条件收敛与重排}
\label{chap:conditional-convergence-rearrangement-solutions}

\begin{solution}{alternating-remainder-sign}
若 \(n\) 为偶数，则 \(s_n\le s\le s_{n+1}=s_n+b_{n+1}\)，故
\[
 0\le s-s_n\le b_{n+1}.
\]
若 \(n\) 为奇数，则
\[
 s_n-b_{n+1}=s_{n+1}\le s\le s_n,
\]
所以 \(-b_{n+1}\le s-s_n\le0\)。两式合起来给出
\(0\le(-1)^n(s-s_n)\le b_{n+1}\)。余项可能为零，例如连续两项相等
的 \(b_{2k-1}=b_{2k}=1/k\) 使偶数部分和恰等于总和零。
\end{solution}

\begin{solution}{alternating-harmonic-error}
Leibniz 估计给出
\[
 \abs{s-s_N}\le\frac1{N+1}.
\]
因此 \(N+1>10^4\)，例如 \(N=10000\)，保证误差小于 \(10^{-4}\)。
这是充分条件；实际误差略小于第一遗漏项。
\end{solution}

\begin{solution}{abel-summation-direct-check}
左端为
\[
 a_2b_2+a_3b_3+a_4b_4.
\]
右端为
\[
 A_4b_4-A_1b_2+A_2(b_2-b_3)+A_3(b_3-b_4).
\]
代入 \(A_2=A_1+a_2\)、\(A_3=A_2+a_3\)、
\(A_4=A_3+a_4\)，\(A_1\) 与各中间项消去，恰剩左端。
\end{solution}

\begin{solution}{dirichlet-geometric-partial-sums}
对 \(x\notin2\pi\Z\)，把三角差公式作望远镜求和，得到
\[
 \sum_{k=1}^n\sin(kx)
 =\frac{\cos(x/2)-\cos((n+1/2)x)}{2\sin(x/2)},
\]
\[
 \sum_{k=1}^n\cos(kx)
 =\frac{\sin((n+1/2)x)-\sin(x/2)}{2\sin(x/2)}.
\]
两个分子的绝对值都不超过 \(2\)，故两个部分和的绝对值都不超过
\(1/|\sin(x/2)|\)。这个界在固定 \(x\notin2\pi\Z\) 时有限，但当 \(x\) 靠近
\(2\pi\Z\) 时变大。
\end{solution}

\begin{solution}{series-sine-over-n}
上一题的正弦部分和公式说明 \(\sum_{k=1}^n\sin(kx)\) 关于 \(n\) 有界。
又 \(1/n\) 单调趋零，Dirichlet 判别给出收敛。
\end{solution}

\begin{solution}{series-abel-monotone-factor}
设 \(A_n=\sum_{k=1}^na_k\to A\)，故 \((A_n)\) 有界。对
\(\sum_{k=m}^na_k(b_k-L)\) 使用 Abel 变换；因 \(b_k-L\) 单调趋零，
尾和由常数乘 \(\abs{b_m-L}\) 控制。故该级数满足 Cauchy 准则。
再加上 \(L\sum a_n\) 即得 \(\sum a_nb_n\) 收敛。
\end{solution}

\begin{solution}{positive-negative-parts-equivalence}
若 \(\sum\abs{a_n}\) 收敛，则
\[
 0\le a_n^\pm\le\abs{a_n},
\]
比较判别给出两者都收敛。反之，若正、负部级数都收敛，则
\(\abs{a_n}=a_n^++a_n^-\)，右侧级数为两个收敛正项级数之和，故绝对值
级数收敛。原和由有限恒等式
\(\sum a_n=\sum a_n^+-\sum a_n^-\) 取极限得到。
\end{solution}

\begin{solution}{subseries-absolute-characterization}
若 \(\sum|a_n|<\infty\)，则任意严格递增指标列 \((n_k)\) 都满足
\[
 \sum_{k=1}^N|a_{n_k}|
 \le\sum_{n=1}^{n_N}|a_n|
 \le\sum_{n=1}^{\infty}|a_n|,
\]
所以对应子级数绝对收敛。

反向假设每个子级数都收敛。全指标列 \(n_k=k\) 给出原级数收敛。若原级数
不绝对收敛，它便条件收敛；正项子级数与负项子级数分别发散到
\(+\infty\) 和 \(-\infty\)，矛盾。因此原级数绝对收敛。

“存在一个收敛子级数”不提供这种控制。例如调和级数发散，但取
\(n_k=2^k\) 得子级数 \(\sum2^{-k}\) 收敛。
\end{solution}

\begin{solution}{conditional-both-signs-infinite}
若正项只有有限多个，则 \(\sum a_n^+\) 收敛。原级数收敛会迫使
\(\sum a_n^-\) 也有有限极限，于是原级数绝对收敛，矛盾。负项有限同理。
\end{solution}

\begin{solution}{absolute-rearrangement-cauchy}
绝对收敛给出：对任意 \(\varepsilon>0\)，存在有限集
\(\{1,\ldots,N\}\)，其外全部项的绝对值总和小于 \(\varepsilon\)。
在任意重排中，当前 \(N\) 项都出现以后，任意后续尾段只由指标大于
\(N\) 的项组成，故尾段绝对值小于 \(\varepsilon\)。重排满足 Cauchy
准则。其部分和与原和之差也只由尚未出现的 \(n>N\) 项构成，故极限相同。
\end{solution}

\begin{solution}{conditional-even-odd-rearrangement}
交错调和级数的正项为 \(1/(2k-1)\)，负项为 \(-1/(2k)\)。
取前 \(2m\) 个正项与前 \(m\) 个负项后的部分和为
\[
 T_m=\sum_{k=1}^{2m}\frac1{2k-1}
      -\frac12\sum_{k=1}^{m}\frac1k
 =H_{4m}-\frac12H_{2m}-\frac12H_m.
\]
用正文引理 \ref{lem:harmonic-log-correction} 的渐近式，得到
\[
 T_m\to\log4-\frac12\log2=\frac32\log2.
\]
每个三项块内部与块末至多相差两项，它们的绝对值和趋于零，故整个重排
也收敛于 \(\frac32\log2\)。原级数的偶数部分和为
\(H_{2m}-H_m\to\log2\)，奇数部分和与其相差 \(1/(2m+1)\to0\)，
所以原和确为 \(\log2\)。
\end{solution}

\begin{solution}{riemann-target-construction-detail}
令正项指标按原序为 \(p(1)<p(2)<\cdots\)，负项指标为
\(q(1)<q(2)<\cdots\)。令 \(r_0=t_0=0\)，空和为零。第 \(j\) 轮中，
取最小的 \(r_j>r_{j-1}\)，使
\[
 \sum_{k=1}^{r_j}a_{p(k)}+\sum_{k=1}^{t_{j-1}}a_{q(k)}>s;
\]
再取最小的 \(t_j>t_{j-1}\)，使
\[
 \sum_{k=1}^{r_j}a_{p(k)}+\sum_{k=1}^{t_j}a_{q(k)}<s.
\]
两种尾和无界，故这两个最小整数存在。按各自原顺序列出本轮新增指标，
得到有限的正项块和负项块。

构造从不重复指标。每阶段至少新增一个对应符号的指标，故
\(r_j,t_j\to\infty\)，所有非零项最终出现；把零项按原序每阶段至多
插入一个，直至穷尽。
因此所得指标列无重复且覆盖 \(\Npos\)，定义了双射。越界误差不超过阶段
最后一项（初始正项阶段可能例外），而通项趋零；块内部分和夹在相邻
越界值之间，所以全部部分和趋于 \(s\)。
\end{solution}

\begin{solution}{riemann-divergent-rearrangements}
趋于 \(+\infty\) 时：第 \(j\) 阶段按序至少取一个尚未使用的正项，继续到部分和超过
\(j\)，再插入一个尚未使用的负项。负项趋零，故阶段末最终仍高于
\(j-1\)。阶段内部先只加正项，故不低于上一阶段末；最后减去负项后的
值仍满足该下界。因此全部部分和趋于无穷。每轮两个符号都使用新项，
所有非零项最终出现；零项每阶段至多插入一个。

上下无界时：第 \(j\) 阶段取正项直到超过 \(j\)，再取负项直到低于
\(-j\)。阶段端点分别趋于 \(+\infty\) 与 \(-\infty\)，所以部分和
上下无界且不收敛。所有正负项按序穷尽，零项依照上一种算法插入。
\end{solution}

\begin{solution}{alternating-not-monotone-counterexample}
令
\[
 b_{2k-1}=\frac1{\sqrt k}+\frac1k,\qquad
 b_{2k}=\frac1{\sqrt k}.
\]
则 \(b_n>0\) 且 \(b_n\to0\)，但每一对交错项之和为
\[
 b_{2k-1}-b_{2k}=\frac1k.
\]
偶数部分和是调和部分和，故级数发散。
\end{solution}

\begin{solution}{dirichlet-boundedness-necessary-example}
取 \(a_n=1\)、\(b_n=1/n\)。则 \(b_n\downarrow0\)，但
\(a_n\) 的部分和为 \(n\)，无界；乘积级数就是调和级数，发散。
\end{solution}

\begin{solution}{conditional-subseries-range}
设正项按原序为 \(p_k>0\)，则 \(p_k\to0\) 且 \(\sum p_k=\infty\)。
给定 \(s>0\)，令尚未补足的差额 \(d_0=s\)。第 \(j\) 阶段先跳过有限
项，直到所有后续 \(p_k<d_{j-1}/4\)，且指标大于上一阶段的末指标。
从此处连续选取正项，到所选块和 \(B_j\) 第一次达到 \(d_{j-1}/2\)。
正项尾和发散保证有限步达到，而最后一项小于 \(d_{j-1}/4\)，故
\[
 \frac{d_{j-1}}2\le B_j<\frac{3d_{j-1}}4.
\]
令 \(d_j=d_{j-1}-B_j\)，则
\(0<d_j\le d_{j-1}/2\)，从而 \(d_j\to0\)。每阶段选择至少一项，
且各阶段指标严格递增，因此所得指标列无限，是原级数的一个子级数。
块末和为 \(s-d_j\to s\)，块内和介于相邻块末和之间，故该子级数收敛
于 \(s\)。保留正的差额，避免了有限步恰好达到 \(s\) 后无项可选的问题。
\end{solution}

\begin{solution}{complex-rearrangement-warning}
复数没有与加法、乘法相容的全序，不能把一个复数项分成“正项”和“负项”
并用越过目标的方式交替逼近。可以分别考察实部和虚部，但同一次重排必须
同时安排两个坐标，两个坐标的要求可能互相制约。因此一维证明不能逐字
照搬。
\end{solution}

\begin{solution}{levy-steinitz-finite-dimensional-note}
取两个条件收敛实级数 \(\sum a_n,\sum b_n\)，形成向量项列
\[
 (a_1,0),(0,b_1),(a_2,0),(0,b_2),\ldots.
\]
给定 \((s,t)\in\R^2\)，分别由 Riemann 重排定理选择排列
\(\pi,\sigma\)，使
\[
 \sum_{k=1}^{\infty}a_{\pi(k)}=s,\qquad
 \sum_{k=1}^{\infty}b_{\sigma(k)}=t.
\]
再按
\[
 (a_{\pi(1)},0),(0,b_{\sigma(1)}),
 (a_{\pi(2)},0),(0,b_{\sigma(2)}),\ldots
\]
交错排列。偶数号向量部分和为
\[
 \left(\sum_{k=1}^{n}a_{\pi(k)},
       \sum_{k=1}^{n}b_{\sigma(k)}\right)\longrightarrow(s,t),
\]
奇数号部分和与相邻偶数号部分和只差一个趋于零的向量项，故也趋于
\((s,t)\)。于是这个向量级数的重排和可以充满整个平面。

若所有向量项都落在一条过原点的直线 \(\R v\) 上，则每个有限部分和以及
每个可能的有限重排极限仍在闭集 \(\R v\) 中，可达到集合不可能充满
\(\R^2\)。本例表明高维重排问题不仅取决于正负号，还取决于条件收敛项
所张成的线性方向；一般结论需要 Lévy--Steinitz 定理，本题不调用该定理。
\end{solution}
